%!TEX root = ../dokumentation.tex

\chapter{Einleitung}

Lorem ipsum dolor sit amet, consetetur sadipscing elitr, sed diam nonumy eirmod tempor invidunt ut labore et dolore magna aliquyam erat, sed diam voluptua. At vero eos et accusam et justo duo dolores et ea rebum. Stet clita kasd gubergren, no sea takimata sanctus est Lorem ipsum dolor sit amet. Lorem ipsum dolor sit amet, consetetur sadipscing elitr, sed diam nonumy eirmod tempor invidunt ut labore et dolore magna aliquyam erat, sed diam voluptua. At vero eos et accusam et justo duo dolores et ea rebum. Stet clita kasd gubergren, no sea takimata sanctus est Lorem ipsum dolor sit amet.

\section{Hintergrund und Motivation}

\begin{figure}[ht]
    \centering
    \includegraphics[height=6cm]{dhbw.png}
    \caption{Exemplarische Abbildung: Das Logo der Dualen Hochschule Baden-Württemberg.}
\end{figure}

\section{Ziele der Arbeit}

Die \ac{DHBW} wurde am 1. März 2009 gegründet \cite{DHBW-wirueberuns}.

\begin{minipage}{\linewidth}
    \begin{lstlisting}[language=Python, caption={Exemplarisches Code-Listing: Einfache for-Schleife mit Konsolenausgabe.}]
    for i in range(10):
        print(i, i * i)
    \end{lstlisting}
\end{minipage}

\begin{table}[ht]
    \centering
    \begin{tabular}{|l|l|}
        \hline
        t {[}s{]} & v {[}m/s{]} \\ \hline
        0         & 0           \\ \hline
        1         & 4           \\ \hline
        2         & 6           \\ \hline
        3         & 7           \\ \hline
    \end{tabular}
    \caption{Exemplarische Tabelle: Messdaten einer Geschwindigkeitsmessung.}
\end{table}

Abbildungen, Tabellen und Code-Listings werden in den jeweiligen Verzeichnissen aufgelistet.